% ch07_u06 —— 原书 p.286--p.295（7.10.1 末段至 7.10.20 证明前半）
% 承接 p.285 末尾 7.10.1 的断句「Set」，本单元自 $f^{!}:=H^{0}Rf^{!}$ 续译；
% 尾部停在 p.295 末行，7.10.20 的证明接 p.296。

\SourcePage{286}

$f^{!}:=H^{0}Rf^{!}:\mathcal{M}(X,\mathcal{O})\to\mathcal{M}(Y,\mathcal{O})$；这是左正合函子。因此，诸范畴 $\mathcal{M}(X,\mathcal{O})$ 连同函子 $f^{!}$，构成概形与拟有限态射的范畴上的纤维范畴。

下面给出 $f^{!}$ 的显式描述。由 Zariski 主定理，任意拟有限态射都是有限态射与开浸入的复合；我们在这两种情形下分别描述 $f^{!}$。若 $f$ 是开浸入（或更一般地，$f$ 平展），则 $f^{!}=f^{*}$。若 $f$ 有限，则 $f^{!}$ 是函子 $f_{*}:\mathcal{M}(Y,\mathcal{O})\to\mathcal{M}(X,\mathcal{O})$ 的右伴随。确切地说，$f_{*}\mathcal{O}_{Y}$ 是有限 $\mathcal{O}_{X}$-代数；函子 $f_{*}$ 把 $\mathcal{M}(Y,\mathcal{O})$ 与这样的范畴等同：其对象是作为 $\mathcal{O}_{X}$-模为拟凝聚的 $f_{*}\mathcal{O}_{Y}$-模。现在设 $M$ 是 $X$ 上的 $\mathcal{O}$-模，则相应的 $f_{*}\mathcal{O}_{Y}$-模 $f_{*}f^{!}M$ 是 $\mathcal{H}om_{\mathcal{O}_{X}}(f_{*}\mathcal{O}_{Y},M)$。特别地，若 $f$ 是闭浸入，则 $f^{!}M\subset M$ 是（概形意义下）支撑于 $Y$ 上的截面构成的子模。

上述图景按如下方式推广到形式概形（或代数空间）的设定。对形式概形 $\hat X$，记 $\mathcal{M}(\hat X,\mathcal{O})$ 为离散拟凝聚 $\mathcal{O}_{\hat X}$-模的范畴\footnote{该范畴是阿贝尔范畴。更一般的设定见 7.11.4。}。例如，设 $\hat X$ 是概形 $V$ 沿其闭子概形 $X$ 的形式完备化，则 $\mathcal{M}(\hat X,\mathcal{O})$ 与 $V$ 上集合论意义下支撑于 $X$ 的 $\mathcal{O}$-模范畴相同。若 $\hat X$ 仿射，则对任意 $M\in\mathcal{M}(\hat X,\mathcal{O})$ 有 $M=\bigcup M_{X'}$；这里 $X'$ 取遍 $\hat X$ 的闭子概形构成的（有向）集合，而 $M_{X'}\in\mathcal{M}(X',\mathcal{O})$ 是概形意义下支撑于 $X'$ 的截面构成子模。拉回函子 $f^{!}$ 以唯一的方式\footnote{我们假设它们与 $f$ 的复合相容。}推广到形式代数空间的拟有限态射的设定。事实上，设 $\hat f:\hat Y\to\hat X$ 是这样的态射。为定义 $\hat f^{!}:\mathcal{M}(\hat X,\mathcal{O})\to\mathcal{M}(\hat Y,\mathcal{O})$，可设 $\hat X$、$\hat Y$ 仿射。此时 $\hat f^{!}M=\bigcup\hat f^{!}|_{Y'}M_{X'}$，其中 $Y'$ 是 $\hat Y$ 的闭子概形，且 $\hat f(Y')\subset X'$。若 $\hat f$ 是 ind-有限\footnote{$:=Y_{\mathrm{red}}\to X_{\mathrm{red}}$ 是有限态射。}，则 $\hat f^{!}$ 的显式描述留给读者。

\SourcePage{287}

\bdp{7.10.2} 对概形或代数空间 $X$，记 $X_{cr}$ 为如下图表构成的范畴：$X\overset{j}{\longleftarrow}S\overset{i}{\hookrightarrow}\hat S$，其中 $j$ 是拟有限态射，$i$ 是仿射概形的闭浸入，且相应的理想 $\mathcal{I}\subset\mathcal{O}_{\hat S}$ 幂零。我们通常把 $X_{cr}$ 的这个对象写作 $(S,\hat S)$，或简写作 $\hat S$。$X_{cr}$ 中的态射 $(S,\hat S)\to(S',\hat S')$ 是概形态射 $\phi:\hat S\to\hat S'$，使得 $\phi(S)\subset S'$，且 $\phi|_{S}:S\to S'$ 是 $X$-概形的态射。

注意，对如上任意 $\phi$，态射 $\phi:\hat S\to\hat S'$ 都拟有限。因此，诸范畴 $\mathcal{M}(\hat S,\mathcal{O})$ 连同拉回函子 $\phi^{!}$，构成 $X_{cr}$ 上的纤维范畴 $\mathcal{M}^{!}(X_{cr},\mathcal{O})$。

有时考虑更大的范畴 $X_{\widehat{cr}}$ 较为方便：它由与上述类似的图表构成，但允许 $\hat S$ 是形式概形（于是 $\mathcal{I}$ 是投射幂零理想，即 $\hat S_{\mathrm{red}}=S_{\mathrm{red}}$）。与上面一样，有 $X_{\widehat{cr}}$ 上的纤维范畴 $\mathcal{M}^{!}(X_{\widehat{cr}},\mathcal{O})$。

\bdp{7.10.3} \emph{定义.} $X$ 上的 $\mathcal{D}$-晶体是 $\mathcal{M}^{!}(X_{cr},\mathcal{O})$ 的笛卡儿截面。$X$ 上的 $\mathcal{D}$-晶体构成 $\mathbb{C}$-范畴 $\mathcal{M}_{\mathcal{D}}(X)$。

确切地说，$\mathcal{D}$-晶体 $M$ 是这样一条规则：对任意 $(S,\hat S)\in X_{cr}$，指定 $\hat S$ 上的 $\mathcal{O}$-模 $M_{\hat S}=M_{(S,\hat S)}$；对态射 $\phi:(S,\hat S)\to(S',\hat S')$，指定等同 $\alpha_{\phi}:M_{\hat S}\xrightarrow{\sim}\phi^{!}M_{\hat S'}$，且这些等同与 $\phi$ 的复合相容。

特别地，若 $\phi$ 是由理想 $\mathcal{I}\subset\mathcal{O}_{\hat S'}$ 定义的闭浸入，则 $M_{\hat S}$ 是 $M_{\hat S'}$ 中被 $\mathcal{I}$ 零化的截面构成子模。

在定义 7.10.3 中可以把 $X_{cr}$ 换成 $X_{\widehat{cr}}$：得到的 $\mathcal{D}$-晶体范畴相同。事实上，对 $(S,\hat S)\in X_{\widehat{cr}}$ 有 $M_{\hat S}=\bigcup M_{(S,\hat S')}$，其中 $\hat S'$ 取遍所有满足 $S\subset\hat S'\subset\hat S$ 的子概形。

\bdp{7.10.4} \emph{变体.} 设 $X_{cr}^{(i)},\ldots,X_{cr}^{(iv)}$ 是 $X_{cr}$ 的完全子范畴，其对象 $(S,\hat S)$ 分别满足下列条件之一（在 (ii)--(iv) 中设 $X$ 是概形）：
\begin{enumerate}[label=(\roman*),leftmargin=2.2em]
\item $S\to X$ 平展。
\item $S\to X$ 是开浸入。
\item （设 $X$ 仿射）$S\xrightarrow{\sim}X$。
\SourcePage{288}
\item $S\to X$ 是局部闭浸入。
\end{enumerate}

记 $\mathcal{M}^{(i)}_{\mathcal{D}}(X),\ldots,\mathcal{M}^{(iv)}_{\mathcal{D}}(X)$ 为 $\mathcal{M}^{!}(X_{cr},\mathcal{O})$ 在相应子范畴 $X_{cr}^{(a)}$ 上的笛卡儿截面构成的范畴。有显然的限制函子 $\mathcal{M}_{\mathcal{D}}(X)\to\mathcal{M}^{(a)}_{\mathcal{D}}(X)$。这些函子是范畴等价，其验证留给读者\footnote{只需注意：若把 $\mathcal{M}_{\mathcal{D}}(X)$ 换成 $\mathcal{M}^{(a)}_{\mathcal{D}}(X)$，7.10.6、7.10.7、7.10.8 连同其证明逐字仍然成立。}。

\emph{注.} 范畴 $X_{cr}^{(ii)}$ 就是 $X$ 的 Grothendieck 晶体 site（的底范畴），故按 Grothendieck 的术语，$\mathcal{D}$-晶体就是纤维范畴 $\mathcal{M}^{!}(X_{cr}^{(ii)},\mathcal{O})$ 的晶体。我们把 $X_{cr}$ 当作基本设定，因为它直接推广到归纳概形的设定（见 7.11.6）。

\bdp{7.10.5} 设 $f:Y\to X$ 是拟有限态射。它给出忠实函子 $Y_{cr}\to X_{cr}$，把 $Y\overset{j}{\longleftarrow}S\overset{i}{\hookrightarrow}\hat S$ 送到 $Y\overset{fj}{\longleftarrow}S\overset{i}{\hookrightarrow}\hat S$。%
% [Tutor Note: 原书 p.288 此处第二个图表仍自 $Y$ 起头，按定义（进而给出 $X_{cr}$ 的对象）应自 $X$ 起头；底本如此，保留原样。]
于是得到相应的「限制」函子 $f^{!}:\mathcal{M}_{\mathcal{D}}(X)\to\mathcal{M}_{\mathcal{D}}(Y)$。它与 $f$ 的复合相容。

特别地，当 $U$ 在 $X$ 上平展时，诸范畴 $\mathcal{M}_{\mathcal{D}}(U)$ 构成小平展 site $X_{\text{ét}}$ 上的纤维范畴，记作 $\mathcal{M}_{\mathcal{D}}(X_{\text{ét}})$。

\bdp{7.10.6} \emph{引理.} $\mathcal{D}$-晶体是平展拓扑下的局部对象，即 $\mathcal{M}_{\mathcal{D}}(X_{\text{ét}})$ 是范畴的层。\hfill$\Box$

\bdp{7.10.7} 下面给出 $\mathcal{D}$-晶体的一个方便的「具体」描述。

设有闭浸入 $X\hookrightarrow V$，其中 $V$ 是形式光滑\footnote{见 7.11.1。}的形式代数空间，且 $X_{\mathrm{red}}=V_{\mathrm{red}}$\footnote{即 $X$ 在 $\mathcal{O}_{V}$ 中的理想投射幂零。}。若 $X$ 仿射，这样的对象总存在：可把 $X$ 嵌入光滑概形 $W$，再取 $V$ 为 $W$ 沿 $X$ 的形式完备化。

对 $n\geq1$，记 $V^{<n>}$ 为 $V^{n}$ 沿对角线 $V\subset V^{n}$（等价地，沿 $X\subset V^{n}$）的形式完备化。投影 $p_{1},p_{2}$：\SourcePage{289}$V^{<2>}\to V$，$p_{12},p_{23},p_{13}:V^{<3>}\to V^{<2>}$ 都是 ind-有限，于是有函子 $p^{!}_{i}:\mathcal{M}(V,\mathcal{O})\to\mathcal{M}(V^{<2>},\mathcal{O})$、$p^{!}_{ij}:\mathcal{M}(V^{<2>},\mathcal{O})\to\mathcal{M}(V^{<3>},\mathcal{O})$。由于 $V$ 形式光滑，这些函子正合。

记 $\mathcal{M}_{\mathcal{D}V}(X)$ 为满足下述条件的对 $(M_{V},\tau)$ 构成的范畴：$M_{V}\in\mathcal{M}(V,\mathcal{O})$，而 $\tau:p^{!}_{1}M_{V}\xrightarrow{\sim}p^{!}_{2}M_{V}$ 是同构，且
\begin{equation}\tag{340}
p^{!}_{23}(\tau)p^{!}_{12}(\tau)=p^{!}_{13}(\tau).
\end{equation}

\bdp{7.10.8} \emph{命题.} 范畴 $\mathcal{M}_{\mathcal{D}}(X)$ 与 $\mathcal{M}_{\mathcal{D}V}(X)$ 典范等价。

\emph{证明.} 所论对象是局部的，故可设 $X$ 仿射。对 $M\in\mathcal{M}_{\mathcal{D}}(X)$，有 $M_{V}=M_{(X,V)}\in\mathcal{M}(V,\mathcal{O})$。由于 $p^{!}_{i}M_{V}=M_{V^{<2>}}$，所得的 $\tau$ 显然满足 (340)。反之，设给定 $(M_{V},\tau)\in\mathcal{M}_{\mathcal{D}V}(X)$；我们来定义相应的 $\mathcal{D}$-晶体 $M$。对 $(S,\hat S)\in X_{cr}$，选取延拓结构态射 $j:S\to X$ 的 $j':\hat S\to V$；这样的 $j'$ 存在，因为 $V$ 形式光滑。考虑 $\mathcal{O}_{\hat S}$-模 $j'^{!}M_{V}$。若 $j'':\hat S\to V$ 是 $j$ 的另一延拓，则有典范同构 $\nu_{j'j''}:j'^{!}M_{V}\xrightarrow{\sim}j''^{!}M_{V}$。确切地说，$(j',j'')$ 把 $\hat S$ 映入 $V^{<2>}$，故 $j'^{!}M_{V}=(j',j'')^{!}p^{!}_{1}M_{V}$；再用 $j''^{!}M_{V}$ 的类似描述，令 $\nu_{j'j''}:=(j',j'')^{!}(\tau)$。由 (340)，这些等同可传递，故 $j'^{!}M_{V}$ 不依赖于 $j'$ 的选取。这就是 $M_{(S,\hat S)}$。$M$ 的结构同构 $\alpha_{\phi}$ 的定义是显然的。\hfill$\Box$

\bdp{7.10.9} \emph{推论.} (i) 对任意 $X$，范畴 $\mathcal{M}_{\mathcal{D}}(X)$ 是阿贝尔范畴。

(ii) 对 $\hat S\in X_{cr}$，函子 $\mathcal{M}_{\mathcal{D}}(X)\to\mathcal{M}(\hat S,\mathcal{O})$，$M\mapsto M_{\hat S}$ 左正合。

(iii) 对拟有限态射 $j:Y\to X$，函子 $j^{!}:\mathcal{M}_{\mathcal{D}}(X)\to\mathcal{M}_{\mathcal{D}}(Y)$ 左正合；若 $j$ 平展，则 $j^{!}$ 正合。

\emph{证明.} 命题 (i) 在 $X$ 仿射时成立。事实上，按 7.10.7 选取 $X\hookrightarrow V$。范畴 $\mathcal{M}_{\mathcal{D}V}(X)$ 是阿贝尔范畴，因为函子 $p^{!}_{i}$、$p^{!}_{ij}$ 正合，故由 7.10.8 得证。

\SourcePage{290}

若 $j:U\to X$ 是仿射概形的平展态射，则函子 $j^{!}:\mathcal{M}_{\mathcal{D}}(X)\to\mathcal{M}_{\mathcal{D}}(U)$ 正合。事实上，设 $U\hookrightarrow V_{U}$ 是 $X\hookrightarrow V$ 的 $U$-局部化（故 $V_{U}$ 在 $V$ 上平展）；则 $j^{!}$ 与平展局部化函子 $\mathcal{M}_{\mathcal{D}V}(X)\to\mathcal{M}_{DV_{U}}(U)$ 一致，而后者显然正合。

于是 (i) 由 7.10.6 推出。其余留给读者。\hfill$\Box$

\bdp{7.10.10} \emph{引理.} 对平展态射 $p:U\to X$，函子 $p^{!}$ 有右伴随函子 $p_{*}:\mathcal{M}_{\mathcal{D}}(U)\to\mathcal{M}_{\mathcal{D}}(X)$。若 $p$ 是开浸入，则 $p^{!}p_{*}$ 是恒等函子。

\emph{证明.} 下面给出 $p_{*}$ 的显式构造。对 $(S,\hat S)\in X_{cr}$，令 $S_{U}:=S\times_{X}U$；设 $\hat p_{S}:\hat S_{U}\to\hat S$ 是平展态射，它到 $S\hookrightarrow\hat S$ 的拉回是投影 $S_{U}\to S$。于是 $(S_{U},\hat S_{U})\in U_{cr}$，并有函子 $X_{cr}\to U_{cr}$，$(S,\hat S)\mapsto(S_{U},\hat S_{U})$。

现在对 $N\in\mathcal{M}_{\mathcal{D}}(U)$ 令 $(p_{*}N)_{\hat S}:=(\hat p_{S})_{\cdot}N_{\hat S_{U}}$。等同 $\alpha_{\phi}$ 来自基变换同构 $\phi^{!}\hat p_{S'}{}_{\cdot}=\hat p_{S}{}_{\cdot}\phi^{!}_{U}$。\hfill$\Box$

现在设 $i:Y\hookrightarrow X$ 是闭浸入，$j:U:=X\setminus Y\hookrightarrow X$ 是补开浸入。记 $\mathcal{M}_{\mathcal{D}}(X)_{Y}$ 为 $\mathcal{M}_{\mathcal{D}}(X)$ 的完全子范畴，由满足 $j^{!}M=0$ 的 $\mathcal{D}$-晶体 $M$ 构成。

\bdp{7.10.11} \emph{引理.} (i) 函子 $i^{!}$ 有左伴随函子 $i_{*}:\mathcal{M}_{\mathcal{D}}(Y)\to\mathcal{M}_{\mathcal{D}}(X)$。

(ii) $i_{*}$ 把 $\mathcal{M}_{\mathcal{D}}(Y)$ 映入 $\mathcal{M}_{\mathcal{D}}(X)_{Y}$，且
\[
i_{*}:\mathcal{M}_{\mathcal{D}}(Y)\to\mathcal{M}_{\mathcal{D}}(X)_{Y},\qquad i^{!}:\mathcal{M}_{\mathcal{D}}(X)_{Y}\to\mathcal{M}_{\mathcal{D}}(Y)
\]
是互逆的范畴等价。

(iii) 设 $p:Z\to X$ 是拟有限态射；令 $Y_{Z}:=Y\times_{Z}X$，于是有 $i_{Z}:Y_{Z}\hookrightarrow Z$ 与 $p_{Y}:Y_{Z}\to Y$。则有函子的典范等同 $p^{!}i_{*}=i_{Y*}p^{!}_{Y}:\mathcal{M}_{\mathcal{D}}(Y)\to\mathcal{M}_{\mathcal{D}}(Z)$。

\emph{证明.} 下面给出 $i_{*}$ 的显式构造。取 $Y$ 上的 $\mathcal{D}$-晶体 $N$。对 $(S,\hat S)\in X_{cr}$ 令 $S_{Y}:=S\times_{X}Y$，于是 $S_{Y}$ 是 $S$ 的闭子概形，\SourcePage{291}从而是 $\hat S$ 的闭子概形。投影 $S_{Y}\to Y$ 拟有限，故 $N$ 给出 $S_{Y}$ 上的 $\mathcal{D}$-晶体。把 $(i_{*}N)_{(S,\hat S)}$ 定义为 $\hat S$ 上相应的 $\mathcal{O}$-模（见 7.10.3）。$i_{*}N$ 的结构同构 $\alpha_{\phi}$ 按显然的方式来自 $N$ 的相应同构。

$i_{*}$ 的伴随性质以及性质 (ii)、(iii) 都是显然的。\hfill$\Box$

\bdp{7.10.12} \emph{命题.} 若 $X$ 光滑，则 $\mathcal{M}_{\mathcal{D}}(X)$ 典范等价于 $X$ 上 $\mathcal{D}$-模的范畴 $\mathcal{M}(X)$。

\emph{证明.} 对 $V=X$ 使用 $\mathcal{M}_{\mathcal{D}}(X)$ 的描述 7.10.7。于是 $\mathcal{D}$-晶体 $M$ 等同于对 $(M_{X},\tau)$，其中 $M_{X}\in\mathcal{M}(X,\mathcal{O})$，而 $\tau:p^{!}_{1}M_{X}\xrightarrow{\sim}p^{!}_{2}M_{X}$ 是 $X^{<2>}$ 上 $\mathcal{O}$-模的同构，且满足 (340)。我们来说明，这样的 $\tau$ 与 $M_{X}$ 上的右 $\mathcal{D}$-模结构是同一回事。

把 $\mathcal{D}_{X}$ 看作 $\mathcal{M}(X^{<2>},\mathcal{O})$ 的对象（经由 $\mathcal{O}_{X}$-双模结构）。有典范同构 $\mathcal{D}_{X}\xrightarrow{\sim}p^{!}_{1}\mathcal{O}_{X}$，它把 $\partial\in\mathcal{D}_{X}$ 与截面 $(f\otimes g\mapsto f\partial(g))\in\mathcal{H}om_{\mathcal{O}_{X}}(\mathcal{O}_{X^{<2>}},\mathcal{O}_{X})=p^{!}_{1}\mathcal{O}_{X}$ 对应。因此 $M_{X}\otimes_{\mathcal{O}_{X}}\mathcal{D}_{X}\xrightarrow{\sim}M_{X}\otimes_{\mathcal{O}_{X}}p^{!}_{1}\mathcal{O}_{X}\xrightarrow{\sim}p^{!}_{1}M_{X}$。于是由伴随性，
\begin{equation}\tag{341}
\operatorname{Hom}(p^{!}_{1}M_{X},p^{!}_{2}M_{X})=\operatorname{Hom}(p_{2\cdot}p^{!}_{1}M_{X},M_{X})=\operatorname{Hom}(M_{X}\otimes\mathcal{D}_{X},M_{X}).
\end{equation}
这里把 $M_{X}\otimes\mathcal{D}_{X}$ 经由 $\mathcal{D}_{X}$ 上的右 $\mathcal{O}$-模结构看作 $\mathcal{O}_{X}$-模。于是 $\tau:p^{!}_{1}M_{X}\to p^{!}_{2}M_{X}$ 等同于态射 $M_{X}\otimes\mathcal{D}_{X}\to M_{X}$。可以验证，对 $\tau$ 的条件恰好意味着这个箭头是 $\mathcal{D}_{X}$ 在 $M_{X}$ 上的右幺作用。关于评注与细节，见下一个注记。\hfill$\Box$

\bdp{7.10.13} \emph{注.} 我们在更一般的设定下讨论 7.10.12 的若干要点。由于 $\mathcal{O}_{X^{<2>}}$ 是 $\mathcal{O}_{X}\otimes_{\mathbb{C}}\mathcal{O}_{X}$ 的完备化，可以把 $\mathcal{M}(X^{<2>},\mathcal{O})$ 的对象看作 $X$ 上某些 $\mathcal{O}_{X}$-双模层，称为 Diff-双模\footnote{[BB93] 中用的是「微分双模」这一术语；细节请参看那里。}。若 $A$、$B$ 是 Diff-双模，则 $A\otimes_{\mathcal{O}_{X}}B$ 也是（故 $\mathcal{M}(X^{<2>},\mathcal{O})$ 是幺半范畴）。注意 $A\otimes_{\mathcal{O}_{X}}B$ 实际上（以显然的方式）是\SourcePage{292}$\mathcal{M}(X^{<3>},\mathcal{O})$ 的对象。

由伴随性，对任意 $C\in\mathcal{M}(X^{<2>},\mathcal{O})$，Diff-双模的态射 $A\otimes_{\mathcal{O}_{X}}B\to C$ 等同于 $\mathcal{M}(X^{<3>},\mathcal{O})$ 中的态射 $A\otimes_{\mathcal{O}_{X}}B\to p^{!}_{13}C$。于是对 Diff-代数\footnote{即 Diff-双模的幺半范畴中的代数。} $A$，其乘法等同于 $\mathcal{M}(X^{<3>},\mathcal{O})$ 中的态射 $m:A\otimes_{\mathcal{O}_{X}}A\to p^{!}_{13}A$（结合律如何用这些术语写出，留给读者）。类似地，对（右）$A$-模 $M_{X}$，可以把 $A$-作用写成 $\mathcal{M}(X^{<2>},\mathcal{O})$ 中的态射 $a:M_{X}\otimes_{\mathcal{O}_{X}}A\to p^{!}_{2}M_{X}$；作用（结合）性质恰好是说，由 $m$ 与 $a$ 得到的 $\mathcal{M}(X^{<3>},\mathcal{O})$ 中两个态射 $M_{X}\otimes_{\mathcal{O}_{X}}A\otimes_{\mathcal{O}_{X}}A\to p^{!}_{3}M$ 相同。现在设 $A=\mathcal{D}_{X}$，或更一般地，$A$ 是 tdo。则 $m:A\otimes_{\mathcal{O}_{X}}A\to p^{!}_{13}A$ 是同构\footnote{这个性质或许刻画了 tdo。}。若 $M_{X}$ 是（可能非幺的）$A$-模，则 $a:M_{X}\otimes_{\mathcal{O}_{X}}A\to p^{!}_{2}M_{X}$ 是同构，当且仅当该模是幺模。

\bdp{7.10.14} 在光滑设定下，把 7.10.5、7.10.10、7.10.11(i) 中的函子 $f^{!}$、$p_{*}$、$i_{*}$ 分别与标准的 $\mathcal{D}$-模函子等同起来，这件事留给读者。

结合 7.10.12 与 7.10.11(ii) 可见：若 $X$ 是任意代数空间，则 $X$ 上的 $\mathcal{D}$-晶体等同于 [Sa91] 意义下 $X$ 上的 $\mathcal{D}$-模\footnote{Saito 更愿意处理解析设定，但他的定义有明显的代数版本（而上述定义有明显的解析版本）。}。

\bdp{7.10.15} 本节其余部分是 tdo 与扭 $\mathcal{D}$-模的晶体设定的概要。先讨论晶体 $\mathcal{O}^{*}$-gerbe。在光滑概形的情形，这样的 gerbe 等同于「tdo 模去线丛的扭」这一概念的平展局部化版本。然后对晶体 $\mathcal{O}^{*}$-gerbe $\mathcal{C}$，定义相应的扭 $\mathcal{D}$-晶体的阿贝尔范畴 $\mathcal{M}_{\mathcal{C}}(X)$。

\bdp{7.10.16} 与前面一样，$X$ 是任意代数空间。范畴 $X_{cr}$ 带有 site 的结构（平展晶体拓扑）：覆盖是态射族\SourcePage{293}$\{(S_{i},\hat S_{i})\to(S,\hat S)\}$，使得 $\{\hat S_{i}\to\hat S\}$ 是 $\hat S$ 的平展覆盖。其上带有环层 $\mathcal{O}_{cr}$，其中 $\mathcal{O}_{cr}(S,\hat S)=\mathcal{O}(\hat S)$。于是有相应的可逆元层 $\mathcal{O}^{*}_{cr}$。

\bdp{7.10.17} \emph{定义.} $X$ 上的晶体 $\mathcal{O}^{*}$-gerbe 是 $X_{cr}$ 上的 $\mathcal{O}^{*}_{cr}$-gerbe\footnote{即 [De-Mi] 术语中 $X_{cr}$ 上带 $\mathcal{O}^{*}_{cr}$ 的 gerbe。}。

确切地说，意思如下。考虑 $X_{cr}$ 上的 Picard 群胚层 $\mathcal{P}ic_{cr}$，其中 $\mathcal{P}ic_{cr}(S,\hat S):=\mathcal{P}ic(\hat S)$（即 $\hat S$ 上线丛的 Picard 群胚）。于是 $X$ 上的晶体 $\mathcal{O}^{*}$-gerbe 是 $X_{cr}$ 上的 $\mathcal{P}ic_{cr}$-torsor $\mathcal{C}$（即 $\mathcal{C}$ 是 $X_{cr}$ 上的纤维范畴，带有 $\mathcal{P}ic_{cr}$ 的作用，使得每个纤维 $\mathcal{C}(\hat S)=\mathcal{C}(S,\hat S)$ 是 $\mathcal{P}ic(\hat S)$-torsor），且在 $X_{cr}$ 上局部地 $\mathcal{C}(S,\hat S)$ 非空。

晶体 $\mathcal{O}^{*}$-gerbe 构成 Picard 2-群胚 $\mathcal{G}_{cr}(X)$。gerbe 的等价类群是 $H^{2}(X_{cr},\mathcal{O}^{*}_{cr})$。对一对 gerbe $\mathcal{C}$、$\mathcal{C}'$，态射 $\phi:\mathcal{C}\to\mathcal{C}'$ 构成 $\mathcal{P}ic(X_{cr})$-torsor。这里 $\mathcal{P}ic(X_{cr})$ 是 $X_{cr}$ 上 $\mathcal{O}^{*}_{cr}$-torsor 的 Picard 群胚\footnote{若 $X$ 光滑，则这样的 torsor 等同于 $X$ 上带平坦联络的线丛。}。

\bdp{7.10.18} \emph{注.} (i) 设 $X_{\text{ét}\,cr}$ 是 $X$ 的小平展晶体 site（作为范畴它等于 7.10.4 中的 $X^{(i)}_{cr}$，拓扑由 $X_{cr}$ 诱导）。$X$ 上的晶体 $\mathcal{O}^{*}$-gerbe 经限制给出 $X_{\text{ét}\,cr}$ 上的 $\mathcal{O}^{*}_{cr}$-gerbe。我们得到 gerbe 的 Picard 2-群胚的等价，其验证留给读者\footnote{我们把 $X_{cr}$ 当作基本设定，因为它直接推广到归纳概形的情形，见 7.11.6。}。

(ii) 我们的 $\mathcal{G}_{cr}(X)$ 是与复形 $\tau_{\leq2}R\Gamma(X_{cr},\mathcal{O}^{*}_{cr})=\tau_{\leq2}R\Gamma(X_{\text{ét}\,cr},\mathcal{O}^{*}_{cr})$ 相联系的 Picard 2-群胚。为计算 $R\Gamma$，考察由 $(\mathcal{O}_{cr}/\mathcal{I}_{cr})(S,\hat S)=\mathcal{O}(S)$ 定义的典范理想 $\mathcal{I}_{cr}\subset\mathcal{O}_{cr}$。有环化拓扑的典范态射 $i:X_{\text{ét}}\to X_{\text{ét}\,cr}$，$i^{-1}(S,\hat S)=S$，而 $\mathcal{I}_{cr}$ 适合短正合列 $0\to\mathcal{I}_{cr}\to\mathcal{O}_{cr}\to i_{\cdot}\mathcal{O}_{X}\to 0$。过渡\SourcePage{294}到可逆元层，得到短正合列
\begin{equation}\tag{342}
0\longrightarrow\mathcal{I}_{cr}\xrightarrow{\exp}\mathcal{O}^{*}_{cr}\longrightarrow i_{\cdot}\mathcal{O}^{*}_{X}\longrightarrow 0,
\end{equation}
其中 $\exp$ 是指数映射（由于每个 $\mathcal{I}_{cr}(S,\hat S)$ 都是幂零理想，我们的 $\exp$ 定义得当）。由于 $R\Gamma(X_{\text{ét}\,cr},i_{\cdot}\mathcal{O}^{*}_{X})=R\Gamma(X_{\text{ét}},\mathcal{O}^{*})$。可以用 (343) 来计算 $R\Gamma(X_{cr},\mathcal{O}^{*}_{cr})$。%
% [Tutor Note: 原书 p.294 此处句末印作句号而非逗号，且所引编号印作 (343)——讨论的其实是本页的 (342)，(343) 在 p.295；底本如此，保留原样。]
例如，由于 $H^{0}(X_{cr},\mathcal{I}_{cr})=0$，群 $H^{0}(X_{cr},\mathcal{O}^{*}_{cr})$ 就是 $X$ 上局部常值可逆函数构成的群 $\mathcal{O}^{*}(X)_{con}$。

(iii) 设 $X$ 光滑。令 $\Omega^{\geq1}_{X}:=(0\to\Omega^{1}_{X}\to\Omega^{2}_{X}\ldots)$。按 Grothendieck，有 $R\Gamma(X_{cr},\mathcal{O}_{cr})=R\Gamma(X,\Omega_{X})$ 与 $R\Gamma(X_{cr},\mathcal{I}_{cr})=R\Gamma(X_{cr},\mathit{Cône}(\mathcal{O}_{cr}\to i_{\cdot}\mathcal{O}_{X})[-1])=R\Gamma(X,\Omega^{\geq1}_{X})$。于是 (342) 给出长上同调序列
\[
\begin{gathered}
0\longrightarrow\mathcal{O}^{*}(X)_{con}\longrightarrow\mathcal{O}^{*}(X)\xrightarrow{\ dlog\ }\Omega^{1cl}(X)\longrightarrow H^{1}(X_{cr},\mathcal{O}^{*}_{cr})\longrightarrow\\
\longrightarrow\mathcal{P}ic(X)\xrightarrow{\ c_{1}\ }H^{2}(X,\Omega^{\geq1}_{X})\longrightarrow H^{2}(X_{cr},\mathcal{O}^{*}_{cr})\longrightarrow Br(X)\longrightarrow0.
\end{gathered}
\]
这里 $H^{1}(X_{cr},\mathcal{O}^{*}_{cr})$ 是 $X$ 上带平坦联络的线丛的同构类群。右端的 0 成立，因为 $H^{2}(X_{\text{ét}},\mathcal{O}^{*})=Br(X)$ 是扭群，而 $H^{3}(X_{cr},\mathcal{I}_{cr})$ 是 $\mathbb{C}$-向量空间。

(iv) 若 $X$ 是概形，则可考虑更弱的拓扑 $X_{Zar\,cr}$（作为范畴它等于 7.10.4 中的 $X^{(ii)}_{cr}$）。得到 $X_{Zar\,cr}$ 上 $\mathcal{O}^{*}_{cr}$-gerbe 的相应 Picard 2-群胚 $\mathcal{G}_{Zar\,cr}(X)$。由平展下降，拉回函子 $\mathcal{G}_{Zar\,cr}(X)\to\mathcal{G}_{cr}(X)$ 是 Picard 2-群胚的全忠实态射；即 $\mathcal{G}_{Zar\,cr}(X)$ 是 Zariski 局部平凡晶体 $\mathcal{O}^{*}$-gerbe 的 2-群胚。容易看出\footnote{参看 7.10.22。}，$\mathcal{C}\in\mathcal{G}_{cr}(X)$ 属于 $\mathcal{G}_{Zar\,cr}(X)$，当且仅当 $X_{\text{ét}}$ 上的 $\mathcal{O}^{*}$-gerbe $i_{\cdot}\mathcal{C}$ 是 Zariski 局部平凡的。例如，若 $X$ 光滑，则 $H^{2}(X_{Zar},\mathcal{O}^{*})=0$，故 $\mathcal{G}_{cr}(X)/\mathcal{G}_{Zar\,cr}(X)=Br(X)$。

\SourcePage{295}

\bdp{7.10.19} 下面给出（适当刚性化的）晶体 $\mathcal{O}^{*}$-gerbe 的一个方便的「具体」描述。

设有如 7.10.7 中的 $X\hookrightarrow V$。对 $\mathcal{C}\in\mathcal{G}_{cr}(X)$ 与 $X$ 在 $V$ 中的无穷小邻域 $X'$，有 $\mathcal{P}ic(X')$-torsor $\mathcal{C}(X')$。令 $\mathcal{C}(V):=\varprojlim\mathcal{C}(X')$（$:=$ $\mathcal{C}$ 在诸 $X'$ 构成的有向集上的笛卡儿截面群胚）；这是 $\mathcal{P}ic(V)$-torsor。

考虑这样的对 $(\mathcal{C},\mathcal{E}_{V})$：$\mathcal{C}\in\mathcal{G}_{cr}(X)$，$\mathcal{E}_{V}\in\mathcal{C}(V)$。这样的对象构成 Picard 群胚 $\mathcal{G}^{V}_{cr}(X)$。确切地说，态射 $(\mathcal{C},\mathcal{E}_{V})\to(\mathcal{C}',\mathcal{E}'_{V})$ 是一对 $(F,\nu)$，其中 $F:\mathcal{C}\to\mathcal{C}'$ 是态射，$\nu:F(\mathcal{E}_{V})\xrightarrow{\sim}\mathcal{E}'_{V}$\footnote{注意这样的对没有对称性，故 $\mathcal{G}^{V}_{cr}(X)$ 是通常的群胚（而 $\mathcal{G}_{cr}(X)$ 是 2-群胚）。}。我们来描述 $\mathcal{G}^{V}_{cr}(X)$。

沿用 7.10.7 的记号。设 $\mathcal{R}$ 是 $V^{<2>}$ 上的线丛，$\beta:p^{*}_{12}\mathcal{R}\otimes p^{*}_{23}\mathcal{R}\xrightarrow{\sim}p^{*}_{13}\mathcal{R}$ 是 $V^{<3>}$ 上线丛的同构，使得下述 $V^{<4>}$ 上线丛同构的图表交换（结合条件）：
\begin{equation}\tag{343}
\begin{tikzcd}
\mathcal{R}_{12}\otimes\mathcal{R}_{23}\otimes\mathcal{R}_{34} \arrow[r] \arrow[d] & \mathcal{R}_{13}\otimes\mathcal{R}_{34} \arrow[d]\\
\mathcal{R}_{12}\otimes\mathcal{R}_{24} \arrow[r] & \mathcal{R}_{14}
\end{tikzcd}
\end{equation}
这里 $\mathcal{R}_{ij}$ 是 $\mathcal{R}$ 经投影 $p_{ij}:V^{<4>}\to V^{<2>}$ 的拉回，而箭头来自 $\beta$。

这样的对 $(\mathcal{R},\beta)$ 构成 Picard 群胚 $G(V)$（关于张量积）。

\bdp{7.10.20} \emph{命题.} Picard 群胚 $\mathcal{G}^{V}_{cr}(X)$ 与 $G(V)$ 典范等价。

\emph{证明.} 对 $(\mathcal{C},\mathcal{E}_{V})\in\mathcal{G}^{V}_{cr}(X)$ 令 $\mathcal{R}:=\mathcal{H}om(p^{*}_{1}\mathcal{E}_{V},p^{*}_{2}\mathcal{E}_{V})\in\mathcal{P}ic(V)$，并把 $\beta$ 定义为复合同构；显然 $(\mathcal{R},\beta)\in G(V)$。于是得到 Picard 群胚的态射 $\mathcal{G}^{V}_{cr}(X)\to G(V)$。

% 尾部接 p.296：7.10.20 的证明在 p.296 上续（「The inverse Morphism assigns to $(R,\beta)$…」），本单元至此停笔。
